\section{Related Work}\label{ch3sec:structure}

This section usually surveys prior research. To write your own literature review, open \texttt{chapters/ch3\_content/sections/2\_rw.tex} and replace the placeholder paragraphs with real citations. Use \verb|\cite{key}| for inline citations, \verb|\citeauthor{key}| for author names, and \verb|\citeyear{key}| for years. See Section~\ref{ch3sec:acronyms-citations} for examples.

Each chapter lives in its own directory under \texttt{chapters/}. The entry point is \texttt{main.tex}, which starts with a \verb|\chapter| command, a \verb|\chaptermark| for the running header, and \verb|\glsresetall| to ensure acronyms are expanded again at the start of the chapter.

\begin{table}[h]
\caption{Chapter structure conventions}\label{ch3tab:demo-comparison}
\centering
\begin{tabular}{l l}
\toprule
\thead{Convention} & \thead{Recommendation} \\
\midrule
Chapter folder & \texttt{chapters/your\_chapter/} \\
Entry file & \texttt{main.tex} \\
Section files & \texttt{sections/section\_name.tex} \\
Table files & \texttt{tables/table\_name.tex} \\
\bottomrule
\end{tabular}
\end{table}

\paragraph{Creating a new chapter.} Copy an existing chapter folder, rename it, and add an import command in \texttt{main.tex}. The command should look like \texttt{\textbackslash import\{chapters/your\_chapter\}\{main.tex\}}. Table~\ref{ch3tab:demo-comparison} summarises the conventions.

\paragraph{Adding sections.} Use the standard \verb|\section{...}| command. Subsections and subsubsections work as usual. If you want every section to start on a new page, keep \verb|\sectionbreak| active in the class file.
